\subsection{What the agents do}
\label{sec:behavior}

The reading so far describes outcomes. This one describes the work behind them, which is the part of
an AI-improvement benchmark that does not depend on which reference one divides by. Every
configuration leaves four things the apparatus records without the agent's cooperation: the directories it created while probing the fast metric, the
turns and shell commands its harness took, the source-only patch it finally submitted, and the audited
trajectory the pipeline segments into separately motivated exploration units. Because every configuration is backed by exactly one patch,
the corpus is a \emph{census} of submitted recipes rather than a sample, and the reading below is a
count rather than an estimate.

\subsubsection{What 272 submitted recipes change}
\label{sec:patches}

The census covers the two complete grids for which we hold the submitted diffs: the $180$ Codex and
$100$ Claude Code configurations. Eight of those submitted an empty patch and are excluded here (they
are the source-extraction failures of \S\ref{sec:failures}), leaving $272$ classified patches; Kimi~K3's
ten are outside the classified corpus. We swept each patch's \emph{added} lines for method terms
grouped into eight families, ordered by how much of the training procedure each family replaces
(Table~\ref{tab:behaviors}). The sweep was written to over-generate: it includes terms expected to be
absent, and two such controls return nothing across all $272$ patches, which is what makes the
positive counts readable. Unchanged upstream code never contributes a match, and a term is counted
once per patch however often it recurs.

\textbf{Across the corpus most patches manage the run --- but that split is a property of the
budget, not of the suite.} Budget, checkpointing or logging changes appear in $253$ of $272$ patches
and an optimization knob in $195$, with selection or averaging in $105$ and a capacity or
adapter change in $73$, while $87$ operate on the objective, $66$ add a source of supervision the
recipe did not have, $23$ substitute the learning rule outright and $21$ touch the
data. Taking those four method-level families together, $122$ of $272$ patches ($44.9\%$) change something about how learning
happens. Of the rest, $141$ change only how the run is managed; the remaining $9$ match no family the
sweep detects at all. \emph{That pooled
ratio should not be read as a general statement about coding agents}: it is an average over regimes
that differ by an order of magnitude. At the two lowest Codex effort levels only $3$ of $53$ patches
reach the method level; at \texttt{xhigh} and \texttt{max} $36$ of $60$ do, and $46$ of Claude
Opus~5's $50$ configurations do at every level that harness exposes. The pooled census describes this
corpus; \S\ref{sec:effort} describes what actually drives it. The families are not exclusive --- the median patch matches
$3.03$ of them on average and edits three files --- so a patch that adds a distillation phase also raises the
learning rate for it.

\textbf{The rare families are task-conditioned rather than thinly spread.} All $23$ learning-rule
substitutions fall on four tasks, and added supervision concentrates on the tasks that expose a
stronger signal to learn from. This is the strongest evidence in the corpus that agents read each
repository rather than apply a template: a template would spread the ambitious families evenly and
it does not. \textbf{And no patch touches the scoring path}: across all $272$, not one modifies a file
whose name matches an evaluator, scorer, grader or judge module, so the provenance boundary of
\S\ref{sec:benchmark} is a measured property of this corpus and not only a design intent.

A per-run audit says the same from the other side, with the caveat that it is
machine-produced and author-reviewed rather than independently adjudicated (\S\ref{sec:frames}). Every
configuration is classified on nine fixed scopes, and on the three that decide whether a result is
attributable to the recipe at all the corpus is unanimous: \textbf{across all $290$ configurations, not
one fetched external data, reached the network, or introduced a model the task did not supply.} Reads of the fixed
dataset and of the starting weights are clean but for a handful the audit marks \emph{suspected} or
\emph{inconclusive} rather than violating, and no configuration is found to have carried trained state
forward from exploration into its submission. On the scopes that govern provenance, then, nothing in this corpus supports a
picture of agents carrying trained state across the visibility boundary.

\begin{table}[t]
\centering
\caption{\textbf{What the 272 classified patches change.} The corpus is the Codex and Claude~Code
grids ($180+100$ configurations) less the eight that submitted an empty patch; Kimi~K3 is outside it.
Upper panel: the eight families, ordered by how much of the training procedure each replaces, with the
task each concentrates on most heavily. Families are not exclusive --- a patch matches $3.03$ of them
on average --- so the shares do not sum to one. Lower panel: how often a patch changes the learning procedure rather than the
run, by system and reasoning effort. \textbf{The rise across effort is visible within Sol and Terra
taken separately}, which is what distinguishes it from one variant's behavior pooled into an
average; Luna is the exception, peaking at 5 of 10. The Claude grid shows no trend because Opus~5
is already at its ceiling at the lowest effort that harness exposes.}
\label{tab:behaviors}
\small
\begin{adjustbox}{max width=\linewidth}
\begin{tabular}{llrrl}
\toprule
& \textbf{Family} & \textbf{$n$} & \textbf{share} & \textbf{Concentrated on} \\
\midrule
\multirow{4}{*}{\rotatebox{90}{\footnotesize run}}
& budget, checkpointing or logging & 253 & 93.0\% & preference opt., code RL, agent RL (28/28) \\
& optimization knobs               & 195 & 71.7\% & unlearning (27/28) \\
& selection and averaging          & 105 & 38.6\% & code RL (27/28), weight averaging (26/28) \\
& capacity and adapters            &  73 & 26.8\% & reward modeling (20/28), preference opt.\ (19/28) \\
\midrule
\multirow{4}{*}{\rotatebox{90}{\footnotesize learning}}
& objective surgery                &  87 & 32.0\% & preference opt., agent RL (15/28) \\
& added supervision                &  66 & 24.3\% & one-shot pruning (14/28) \\
& learning-rule substitution       &  23 &  8.5\% & agent RL (15/28), aesthetic FT (6/28) \\
& data intervention                &  21 &  7.7\% & code RL, agent RL (5/28) \\
\midrule
& \textbf{any of the four learning families} & \textbf{122} & \textbf{44.9\%} & \\
\bottomrule
\end{tabular}
\end{adjustbox}

\par\vspace{6pt}
\begin{adjustbox}{max width=\linewidth}
\begin{tabular}{lcccccc}
\toprule
\multicolumn{7}{l}{\emph{Configurations changing the learning procedure, out of 10 per cell}} \\
\textbf{System} & \texttt{none} & \texttt{low} & \texttt{medium} & \texttt{high} & \texttt{xhigh} & \texttt{max} \\
\midrule
Claude Opus 5   & --- & 10 &  8 & 9 & 9 & 10 \\
Claude Sonnet 5 & --- &  6 &  3 & 5 & 5 &  4 \\
GPT-5.6 Sol     &  1  &  1 &  3 & 4 & 7 &  7 \\
GPT-5.6 Terra   &  0  &  0 &  1 & 3 & 6 &  7 \\
GPT-5.6 Luna    &  1  &  0 &  1 & 2 & 5 &  4 \\
\midrule
\emph{GPT grid, out of 30} & \emph{2} & \emph{1} & \emph{5} & \emph{9} & \emph{18} & \emph{18} \\
\bottomrule
\end{tabular}
\end{adjustbox}
\par\vspace{2pt}
{\footnotesize Classification is mechanical: a 42-term regular-expression sweep over each patch's
added lines only, so unchanged upstream code never contributes a match, counted once per
trajectory however often a term recurs. Two terms included as controls return nothing across all
$272$ patches, and no patch modifies a file whose name matches an evaluator, scorer, grader or
judge module; Appendix~\ref{app:behaviors} gives the classification method in full. Claude Code exposes no \texttt{none} effort level. Task names are
abbreviated as in Table~\ref{tab:tasks}.}
\end{table}


\subsubsection{Reasoning effort buys exploration and ambition, not score}
\label{sec:effort}

Holding the task suite and the harness fixed and sweeping the Codex reasoning-effort budget isolates
what effort actually changes (Figure~\ref{fig:effort}). It changes a great deal, and almost none of
it is the result.

\textbf{First, effort buys persistence.} The single largest effect is not on any measure of quality
but on whether the agent keeps working at all: the median configuration at \texttt{none} submits with
$88\%$ of its four-hour window unspent --- it stops after roughly half an hour --- while at
\texttt{max} it uses $96\%$ of the window. Budget consumed rises monotonically across the six levels
($12\%$ to $96\%$), and over the whole suite $61$ of $290$ configurations hand in with more than half
the window unused.

\textbf{Second, effort buys elaboration --- but barely more ideas.} Over the same sweep the median
configuration goes from $4$ to $16.5$ evaluation directories left behind ($4.1\times$), from $153$ to
$733$ agent turns ($4.8\times$), and from a patch that adds $11.5$ lines to one that adds $228$
($19.8\times$), while exploration cost rises $26$-fold and output tokens about ten-fold. Yet the number
of \emph{distinct} experiments barely moves. Reading the audited trajectory record, which segments each
run into separately motivated exploration units, the median Codex configuration runs $5$ of them at
\texttt{none} and $7$ at \texttt{max} --- a factor of $1.4$ against the $4.1$ and $19.8$ above. A
twenty-fold larger patch is not twenty times more hypotheses; it is a comparable number of hypotheses
carrying far more machinery.

\textbf{Ambition scales with it.} Reading the four method-level families as one class gives a single
number per cell, and it rises from $2$ of the $26$ non-empty Codex patches at \texttt{none} to $18$ of $30$ at
\texttt{xhigh} and \texttt{max} --- a $7.8$-fold rise on matched denominators, visible inside GPT-5.6 Sol and
Terra separately and not only in the pool (Table~\ref{tab:behaviors}, lower panel). More reasoning budget does not merely make an agent try more configurations of the same
kind; it makes it attempt a different kind of change.

\textbf{And the score does not move.} Across the sweep the win rate over the
shipped recipe stays inside a narrow band --- $57\%$ at \texttt{none}, a peak of $68\%$ at \texttt{medium}, a trough of
$52\%$ at \texttt{high}, and $65\%$ at the top --- with no trend, and the formal replay is fixed at
$\sim$$11.5$ hours at every level so the compute behind each artifact is constant. Task completion
does rise, from $24/30$ to $29/30$, but for an extraction reason rather than a capability one
(\S\ref{sec:failures}). \textbf{Eight times the time at the desk, nineteen-fold more edited lines and nearly eight times as many
method-level attempts buy nothing measurable in the artifact.} Nor does the persistence pay off on its
own terms: configurations that consume more than three quarters of their window beat the recipe
$62\%$ of the time and those that hand in with over a quarter unspent beat it $64\%$. This is the paper's sharpest behavioural finding,
and it is not a statement about one model: it holds across all three GPT-5.6 variants.

\begin{figure}[t]
\centering
\includegraphics[width=\linewidth]{fig_effort}
\caption{\textbf{What reasoning effort buys, on the Codex grid (30 configurations per level).} Every
quantity is a median shown relative to the lowest effort. As the reasoning budget rises from
\texttt{none} to \texttt{max}, the agent's exploration cost climbs about $26\times$, the number of
probe experiments it runs about $4\times$, and the lines its patch adds about $20\times$ --- while the
win rate over the shipped recipe (bold) never leaves a narrow band around its starting value and
peaks at \texttt{medium}, not at the top. Task completion also rises over the same sweep
($80\%\!\rightarrow\!97\%$), but for an extraction reason (\S\ref{sec:failures}) rather than
capability, and formal retraining is fixed at $\sim$$11.5$ hours per level. Effort buys activity and
ambition, not a better result.}
\label{fig:effort}
\end{figure}

\subsubsection{The harness shapes the search}
\label{sec:footprints}

The two harnesses produce recognisably different work from comparable models, which is the clearest
reason a cross-harness comparison mixes model with scaffold (\S\ref{sec:standings}).

\textbf{Claude Code starts where Codex ends.} Claude Opus~5 reaches the method level in $46$ of its
$50$ configurations, and it is already at $10$ of $10$ at \texttt{low}, the lowest effort that
harness exposes --- the ceiling the Codex grid needs its top two effort levels to approach. Claude
Sonnet~5 sits at $23$ of $50$ and shows no trend in effort either. Across the Codex grid the same
count is $23$ of $60$ for Sol, $17$ for Terra and $13$ for Luna.

\textbf{The gap holds on every task, and narrows only where the task hands out a method.} Comparing
the two harnesses task by task, Claude Code intervenes at the method level more often on \emph{all
ten}, by margins from $18$ to $54$ points. The two narrowest gaps sit at opposite ends of the range and are narrow for
opposite reasons: on graph generation both stay low ($30\%$ against $12\%$), and on agentic RL both run
high ($80\%$ against $61\%$) because the task itself supplies a method-level route
(\S\ref{sec:convergent}). A harness therefore sets a system's default level of intervention, and the
task pulls Codex up to Claude Code's level only when the method is lying in the repository waiting to
be picked up.

\textbf{And the four hours are spent differently.} Counting shell commands across all $290$
exploration logs, Codex spends its window on sheer volume of short actions: its median run issues $48$
commands at \texttt{none} and $507$ at \texttt{max}, a tenfold climb driven entirely by the effort
setting. A Claude Code run issues between $128$ and $170$ at every level, with no
trend, and Kimi~K3 issues $77$. The split is the scaffold's rather than the model's --- per-model
medians run $137$ and $168$ on Claude Code against $128$, $238$ and $321$ on Codex, and only the latter
move with effort. The same four hours buy many cheap actions on one harness and fewer, heavier ones on
the other, which is the mechanism behind the breadth-versus-depth signature below.

\textbf{Breadth versus depth.} The Codex systems run many short probes and submit compact patches:
GPT-5.6 Sol leaves a median $16.5$ evaluation directories, the most of any system, but its median
patch adds $130$ lines, and Terra's and Luna's add $33$ and $25$. The Claude Code systems run fewer,
heavier iterations and submit far larger patches: Claude Opus~5's median patch adds $980$ lines from
a median $14$ evaluation directories. The grids exaggerate the gap: pooled over every level it is roughly eight
times Sol's median and thirty times Terra's, and over the five levels both harnesses expose, five and
twenty times. Kimi~K3, on the
same harness, adds $271$. A larger patch is not a better one --- Sonnet~5 adds $147$ lines and sits fourth
in the standings while Opus~5 leads --- but the gap is consistent, and it says the scaffold decides
\emph{how} a system searches, not only how well.

\subsubsection{Where a system locates the problem}
\label{sec:frames}

The counts so far say how much a system intervenes. The audited trajectories say something the counts
cannot: \emph{where} it decides the problem lies. Each run is segmented into separately motivated exploration units,
$1{,}607$ of them across the corpus, and the first one or two carry the diagnosis everything downstream
serves. That segmentation, and the boundary audit of \S\ref{sec:patches},
come from a language-model pipeline reading each exploration log against a fixed template and reviewed
by the authors against those logs; with no inter-annotator agreement statistic, the unit counts and
frame assignments here are a reviewed machine reading rather than adjudicated labels. The patch census
and the scores below do not depend on them. Sixty-four per cent of runs open with a baseline or diagnostic unit and
$51\%$ close with a delivery unit --- wall-clock stopping, checkpoint retention, export --- so the
canonical shape of a four-hour run is diagnose, intervene, make it survive the replay, with $16\%$ of
all units spent on that last phase alone.

What differs is the diagnosis. Diffusion fine-tuning is the clearest case, because all three families
attack the same repository and reach different conclusions in their first unit.

\begin{itemize}[leftmargin=*,itemsep=2pt,topsep=3pt]
\item \textbf{Claude Opus~5 locates the problem in the update rule.} Its opening units are literally
``Can a differentiable reward replace DDPO/PPO?'' and ``Can direct reward gradients replace the
high-variance DDPO update?''. The hypothesis arrives before any baseline has been shown to
underperform: the agent reads the shipped code, judges the score-function estimator of DDPO
\citep{black2023ddpo} too noisy, and spends its remaining units tuning the replacement --- adapter
capacity, learning rate, truncation depth, safeguards for unattended training. It does this in
\emph{all five} of its configurations. The move it converges on is the one the literature calls direct
reward fine-tuning, in the truncated form that back-propagates through only the last $K$ steps of
sampling \citep{clark2024draft} --- and it inherits that method's documented failure mode, in which the
fine-tuned model collapses onto a single high-reward image \citep{clark2024draft,black2023ddpo}. No
agent cites this literature; each arrives at the method from the repository.
\item \textbf{Claude Sonnet~5, on the same harness, locates it in the schedule.} Its opening units
ask whether the baseline ``is limited by its short schedule rather than its per-epoch recipe'' and
proceed to longer unchanged runs, fused attention kernels, sparse saves and export reliability. It
never questions the update rule in any of its five configurations.
\item \textbf{The Codex systems locate it in the hyperparameters and the train/evaluate mismatch.}
GPT-5.6 Sol establishes a matched baseline, then works through trajectory length, PPO trust region
against optimizer step size, update frequency and advantage normalization, and matches the
evaluator's 20-step sampler. Only at \texttt{max} effort, and only at its fifth unit, does it reach
``directly refine LoRA through a short differentiable aesthetic path''.
\end{itemize}

These are three coherent research programmes, not three levels of competence, and they are chosen
early and from reading rather than from evidence. They also account for the artifacts: Opus~5's five
runs score $8.42$ to $17.67$, Sonnet~5's five sit between $4.49$ and $6.15$ around a recipe at $5.53$,
and Sol crosses $8$ only in the one configuration where its frame finally shifts. \textbf{On this task
a system's default framing separates the artifacts more cleanly than its effort setting does.} We
state that as an observation about the two tasks where the framings diverge visibly --- diffusion
fine-tuning here and agentic RL in \S\ref{sec:convergent} --- and not as a suite-wide effect: on the
remaining eight the diagnoses are closer together, one run per cell gives no noise floor, and we have
not tested whether framing predicts the score where it is not this legible.

\subsubsection{When the frame shifts, and what that is worth}
\label{sec:convergent}

Agentic RL supplies the natural experiment, because the repository ships a breadth-first optimal
Sokoban solver in its own utilities --- on no path the evaluator uses, but readable --- so an
alternative frame is available to any agent that looks. Fourteen of the $29$ submitted patches make some use of it,
across every system but Kimi~K3. The decisive fact is that finding it is not what separates them:
sweeping the exploration logs rather than the patches, \textbf{$27$ of the $29$ runs encounter the
solver, and almost all within the first tenth of their trajectory} --- Claude Opus~5 reads it $18$ to
$34$ lines into logs $773$ to $1{,}017$ lines long, at every effort level, including the two where it
never appears in the patch. Nearly every system saw the same object. \textbf{What separates the
artifacts is what each decided the object was for}, and that decision is legible in when it is taken:

\begin{itemize}[leftmargin=*,itemsep=2pt,topsep=3pt]
\item \textbf{Pivot in the first unit ($7$ configurations): every one that scored reached $0.998$ or
above.} These runs find the solver \emph{during cold-start diagnosis}, not after a failed training
attempt --- Claude Opus~5's opening units are ``Cold-start diagnosis, solver feasibility, and input
format'' and ``Establishing the weak baseline and an available supervised signal'', and the
substitution follows in the next unit.
\item \textbf{Pivot in the second unit ($5$): scores land between $0.27$ and $0.90$}, the frame having
shifted after one round of evidence rather than before it.
\item \textbf{Pivot late or never ($17$, of which two reach the solver only after several failed
repairs): none of the fourteen that scored exceeds $0.27$.} These runs stay inside the shipped
method and try to repair it --- potential-based reward shaping, KL anchors against format drift,
hard-case curricula. The repairs are competent and they do not work.
\end{itemize}

Timing is what the record makes visible, not the mechanism behind it. On this
task the decision and its content are one event --- every patch that builds the training signal on the
solver scores at or above $0.8945$, every patch that keeps the shipped update rule at or below
$0.2656$, nothing between --- and no run decides late and still submits the substitution, so
\emph{when} and \emph{what} cannot be separated here. The negative is what the corpus establishes:
discovery is not the discriminator, so the difference is one of judgement rather than of search.

\textbf{What separates them is the role the agent assigns to the same object.} Opus~5's four
solver-reading runs use it three different ways. At \texttt{low} and \texttt{medium} it becomes a
source of \emph{training labels}: the agent writes a generator whose docstring reads ``build
supervised multi-turn Sokoban trajectories from the in-container solver \dots\ plans come from the
RAGEN BFS shortest-path solver'', and the shipped update rule is gone. At \texttt{xhigh} it becomes a
\emph{shaping signal} for the update rule that stays: the agent builds an exact solve-distance table
and annotates it in its own source as ``training-time only \dots\ it changes the \emph{reward} the
trainer sees, not the environment's success flag or its observations''. At \texttt{max} it becomes an
\emph{instrument}: a prober that establishes ``all 256 public-bank boards solvable in $\le 9$ moves
\dots\ the 10-action cap is not the ceiling'', after which the run returns to repairing the
reinforcement-learning recipe. Three uses, one artifact, scores of $1.000$, $0.211$ and $0.234$.

Reasoning effort moves this dial monotonically. On the Codex grid it shows as
attention --- solver mentions per log climb from $1$ at \texttt{none} to $29$ at \texttt{max} for
GPT-5.6 Sol, and from $1$ to $13$ and $3$ to $16$ for Terra and Luna --- while on Claude Code, where
engagement is already high at every level, it shifts the object's role away from replacement and
toward support. In neither family does an agent explicitly weigh the substitution and reject it; the
record shows something quieter, an agent that has read the same code arriving at a more conservative
use of it as it deliberates longer.

This rests on the one task that supplies a legible alternative route and so makes
the choice observable at all, and we report it as a mechanism we can document rather than an effect we
have estimated. On this task effort did not make a system more willing to abandon the method it was
given; if anything it made a strong system more committed to it, and more inventive about supporting
it. What effort changes is how thoroughly the chosen frame gets executed, which is what
\S\ref{sec:effort} finds across the whole suite.

